% !TeX root = ../../../main.tex
\section{Astronomy: Separate the Sky, the Data Center, and the Sensor}\label{sec:astro}

Astronomy contains at least three distinct quantum opportunities, and combining them creates confusion.

\subsection{Classical astronomical data analysis}

Telescope images, catalogs, and interferometer time series are classical after detection. Quantum algorithms may help with search, inference, linear algebra, sampling, or optimization, but the access and output gates become severe. A quantum image representation that requires loading every pixel and reconstructing every output pixel has likely moved work rather than removed it. Radio-astronomy studies have therefore focused explicitly on comparing quantum encodings and toy reconstruction pipelines, which is the correct early question \citep{brunet2023}.

Gravitational-wave matched filtering is a cleaner mechanism example. After detection, the data are classical, and the physically correct statistic is noise-weighted: given detector data $d$ with one-sided noise power spectral density (PSD) $S_n(f)$---the frequency-resolved variance of the detector noise---and a bank of $M$ waveform templates $h_j$, the inner product is
\begin{equation}
\langle d,h_j\rangle
=4\,\mathrm{Re}\!\int_{0}^{\infty}
\frac{\tilde{d}(f)\,\tilde{h}_j^{*}(f)}{S_n(f)}\,df,
\qquad
\rho_j=\frac{\langle d,h_j\rangle}{\sqrt{\langle h_j,h_j\rangle}},
\label{eq:gw-snr}
\end{equation}
where tildes denote Fourier transforms; templates whose signal-to-noise ratio $\rho_j$ exceeds a threshold are candidate detections. Any quantum proposal must compute \cref{eq:gw-snr}, including the $1/S_n(f)$ weighting and the PSD's own estimation from data, not an unweighted correlation. A reversible predicate $\chi(j)=[\rho_j\geq\rho_*]$ supports Grover-style search, yielding an ideal quadratic dependence on the number of templates \citep{gao2022}. Quantum Monte Carlo and amplitude amplification variants shift the arithmetic and qubit tradeoffs \citep{miyamoto2022}. The decisive questions are whether waveform generation and noise-weighted inner products can be computed coherently, whether data loading is amortized, how many matches exist, and whether the quantum method beats FFT-based and reduced-order classical pipelines at the same false-alarm and false-dismissal rates.

\subsection{Quantum-native cosmological and astrophysical physics}

Some astronomical questions reduce to quantum field theory or many-body physics rather than telescope-data processing: early-universe dynamics, vacuum decay, dense nuclear matter, neutrino transport, axion or dark-sector models, and spectra used to interpret observations. These can inherit the simulation mechanisms and error budgets of subatomic physics. The quantum computer does not simulate an entire galaxy by placing every star in superposition. It may calculate a response function, reaction rate, equation-of-state contribution, or nonperturbative field dynamic that becomes one term inside a classical astrophysical model.

This suggests a productive hybrid hierarchy. The QPU estimates a compact quantity from a quantum many-body model. High-performance classical software propagates that quantity through stellar, galactic, or cosmological scales. Statistical inference then compares the macroscopic prediction with observations. The quantum kernel is deep in the model, not sitting at the telescope's file server.

\subsection{Quantum sensing and networks}

Quantum-enhanced telescopes, photon sorting, entanglement-assisted interferometry, atomic clocks, and quantum detectors are not merely faster algorithms. They change the measurement before classical data exist. This may be the cleaner route to astronomical advantage because the input is quantum-native light. Quantum hypothesis testing and Fisher-information bounds can ask how many photons are fundamentally required to detect a faint companion near a bright star. The output is compact---detection, localization, or parameter precision---and no artificial QRAM is needed.

Quantum sensing must still be separated from quantum computing in claims and budgets---and it sits \emph{outside} the processor-centric framework of this paper. The problem contract of \cref{sec:contract} governs computation on a programmable device; a sensing proposal needs a different resource contract in which loss, dark counts, mode mismatch, baseline distribution, quantum memory lifetime, and network entanglement rates replace gate depth as the controlling resources. NASA's quantum research program appropriately spans algorithms, hardware evaluation, and mission problems rather than assuming one quantum modality fits all astronomy \citep{nasaquail2026}.

\begin{thesisbox}[A provisional assessment for astronomy]
The following ranking is the author's provisional judgment, not a community consensus. Criteria: maturity of the underlying mechanism, severity of the input/output interface, and strength of the classical incumbent, assessed on a horizon of roughly one hardware generation. Under those criteria: (1) quantum-enhanced measurement of quantum light, because the input is quantum-native and no loading problem exists; (2) compact quantum subroutines for native many-body physics inside classical astrophysical models, because the mechanism is Hamiltonian simulation with a compact output; (3) highly structured search or inference with generative data access, where the oracle can be computed rather than loaded. Bulk processing of already classical astronomical archives ranks last: it is the most exposed to loading costs, readout limits, and dequantization. The ranking should be revisited as hardware and classical baselines move.
\end{thesisbox}
